\editionPage[\pdfbookmark[0]{Deux bilans annuels}{bm-bilans-guide-parent}]{bilans-guide}{Deux bilans pour l’année}{Sujets originaux · Relier les chapitres}{Durée conseillée : 1 h 30 · Barème d’entraînement sur 20}
\begin{blocEdition}{bilans-guide}

Ces évaluations sont des créations du livre pour le CM2. Elles n’ont pas
de statut d’examen officiel. Le temps et le barème sont des repères pour
une séance autonome ; ils s’adaptent librement à la classe ou à la maison.
\editionEtape{Automatismes : environ 20 minutes}
Douze questions courtes, sans calculatrice ni brouillon long : 6 points,
soit 0,5 point par réponse. Écrire les étapes utiles, même brèves.
\editionEtape{Problèmes : environ 1 h 10}
Quatre problèmes indépendants : 14 points. Chaque problème vaut
3,5 points : trois questions à 1 point et 0,5 point pour la clarté de la
rédaction, les unités et les contrôles.
\editionEtape{Travailler après l’évaluation}
Comparer les démarches, relever les erreurs par notion, puis refaire
les questions utiles sans correction. Une démarche peut être juste
même quand le résultat final est faux : repérer l’étape à reprendre.
\editionEtape{Préparer une autre séance}
Revenir au cours et à deux exercices du chapitre concerné. Le second
bilan peut être fait plus tard pour vérifier que la méthode est acquise.

\end{blocEdition}
\editionPage{bilan-A-auto}{Bilan A}{Nombres, calculs et fractions}{Partie 1 · 20 min · Sans calculatrice · 6 points}
\begin{blocEdition}{bilan-A-auto}
\begin{questions}
\item Écrire en chiffres : « trois-millions-quarante-mille-six ».
\item Calculer de tête $623-98$.
\item Calculer $40\times700$.
\item Calculer $(15+25)\times(9-4)$.
\item Donner tous les diviseurs de 20.
\item Calculer le quotient décimal exact de $13:4$.
\item Écrire $\frac{17}5$ avec un entier et une fraction.
\item Calculer $\frac34+\frac18$.
\item Calculer les $\frac23$ de 27.
\item Comparer $5{,}09$ et $5{,}1$.
\item Calculer $4{,}56\times100$.
\item Calculer la moitié de $9{,}4$.\end{questions}
\end{blocEdition}
\editionPage{bilan-A-ex-1}{Bilan A · Problème 1}{Le concert de la région}{Partie 2 · 1 h 10 pour les quatre problèmes · 3,5 points}
\begin{blocEdition}{bilan-A-ex-1}
\begin{questions}
\item Un concert a réuni $18\,750$ spectateurs le samedi et $16\,480$ le dimanche. Combien de spectateurs en tout ?
\item Un billet coûte 24 euros. Calculer la recette du samedi.
\item La salle peut accueillir $20\,000$ personnes par soir. Combien de places sont restées vides sur les deux soirs ?
\end{questions}\par\vspace{5mm}{\otLight\fontsize{8}{10}\selectfont\color{otMuted}Recherche et rédaction}\par\par\vspace{7mm}{\color{otRule!65}\hrule height .25pt}\par\vspace{7mm}{\color{otRule!65}\hrule height .25pt}\par\vspace{7mm}{\color{otRule!65}\hrule height .25pt}\par\vspace{7mm}{\color{otRule!65}\hrule height .25pt}\par\vspace{7mm}{\color{otRule!65}\hrule height .25pt}\par\vspace{7mm}{\color{otRule!65}\hrule height .25pt}
\end{blocEdition}
\editionPage{bilan-A-ex-2}{Bilan A · Problème 2}{Les sachets de bonbons}{Partie 2 · 1 h 10 pour les quatre problèmes · 3,5 points}
\begin{blocEdition}{bilan-A-ex-2}
\begin{questions}
\item Un confiseur répartit $1\,000$ bonbons en sachets de 8. Combien de sachets complets prépare-t-il ? Reste-t-il des bonbons ?
\item Un sachet est vendu $2{,}50$ euros. Quelle somme rapportent tous les sachets ?
\item Il partage ensuite 45 euros de pourboire entre 4 vendeurs, au centime près. Combien reçoit chacun ?
\end{questions}\par\vspace{5mm}{\otLight\fontsize{8}{10}\selectfont\color{otMuted}Recherche et rédaction}\par\par\vspace{7mm}{\color{otRule!65}\hrule height .25pt}\par\vspace{7mm}{\color{otRule!65}\hrule height .25pt}\par\vspace{7mm}{\color{otRule!65}\hrule height .25pt}\par\vspace{7mm}{\color{otRule!65}\hrule height .25pt}\par\vspace{7mm}{\color{otRule!65}\hrule height .25pt}\par\vspace{7mm}{\color{otRule!65}\hrule height .25pt}
\end{blocEdition}
\editionPage{bilan-A-ex-3}{Bilan A · Problème 3}{La pizza partagée}{Partie 2 · 1 h 10 pour les quatre problèmes · 3,5 points}
\begin{blocEdition}{bilan-A-ex-3}
\begin{questions}
\item Lina mange $\frac38$ d’une pizza et Tom $\frac14$. Quelle fraction ont-ils mangée ensemble ?
\item Quelle fraction de la pizza reste-t-il ?
\item La pizza entière pesait 480 g. Quelle masse de pizza reste-t-il ?
\end{questions}\par\vspace{5mm}{\otLight\fontsize{8}{10}\selectfont\color{otMuted}Recherche et rédaction}\par\par\vspace{7mm}{\color{otRule!65}\hrule height .25pt}\par\vspace{7mm}{\color{otRule!65}\hrule height .25pt}\par\vspace{7mm}{\color{otRule!65}\hrule height .25pt}\par\vspace{7mm}{\color{otRule!65}\hrule height .25pt}\par\vspace{7mm}{\color{otRule!65}\hrule height .25pt}\par\vspace{7mm}{\color{otRule!65}\hrule height .25pt}
\end{blocEdition}
\editionPage{bilan-A-ex-4}{Bilan A · Problème 4}{Le programme mystère}{Partie 2 · 1 h 10 pour les quatre problèmes · 3,5 points}
\begin{blocEdition}{bilan-A-ex-4}
\begin{questions}
\item Exécuter le programme « choisir un nombre ; multiplier par 3 ; ajouter 6 ; multiplier par 2 » avec le nombre 4.
\item Quel nombre faut-il choisir pour obtenir 48 ?
\item Écrire le résultat du programme pour un nombre $N$ quelconque, avec des parenthèses.
\end{questions}\par\vspace{5mm}{\otLight\fontsize{8}{10}\selectfont\color{otMuted}Recherche et rédaction}\par\par\vspace{7mm}{\color{otRule!65}\hrule height .25pt}\par\vspace{7mm}{\color{otRule!65}\hrule height .25pt}\par\vspace{7mm}{\color{otRule!65}\hrule height .25pt}\par\vspace{7mm}{\color{otRule!65}\hrule height .25pt}\par\vspace{7mm}{\color{otRule!65}\hrule height .25pt}\par\vspace{7mm}{\color{otRule!65}\hrule height .25pt}
\end{blocEdition}
\editionPage{bilan-B-auto}{Bilan B}{Grandeurs, géométrie et données}{Partie 1 · 20 min · Sans calculatrice · 6 points}
\begin{blocEdition}{bilan-B-auto}
\begin{questions}
\item Convertir $2{,}5$ km en mètres.
\item Convertir 750 g en kilogrammes.
\item Convertir $3\ \mathrm{dm}^2$ en $\mathrm{cm}^2$.
\item Aire d’un rectangle de 9 cm sur 6 cm.
\item Convertir 2 min 30 s en secondes.
\item Combien mesure la moitié d’un angle droit ?
\item Combien de côtés a un hexagone ?
\item Combien d’arêtes a un pavé ?
\item Quatre cahiers coûtent 6 euros. Combien en coûtent 12 ?
\item Un dé est lancé. Chances d’obtenir 5 ou 6 ?
\item Un rayon mesure $3{,}5$ cm. Que mesure le diamètre ?
\item Donner le périmètre d’un carré de côté $4{,}5$ cm.\end{questions}
\end{blocEdition}
\editionPage{bilan-B-ex-1}{Bilan B · Problème 1}{La rénovation de la chambre}{Partie 2 · 1 h 10 pour les quatre problèmes · 3,5 points}
\begin{blocEdition}{bilan-B-ex-1}
\begin{questions}
\item Une chambre rectangulaire mesure $4{,}5$ m sur $3{,}2$ m. Calculer son aire.
\item Le parquet coûte 25 euros le mètre carré. Quel est le prix du parquet ?
\item On pose une plinthe sur tout le tour, sauf la porte de $0{,}8$ m. Quelle longueur de plinthe faut-il ?
\end{questions}\par\vspace{5mm}{\otLight\fontsize{8}{10}\selectfont\color{otMuted}Recherche et rédaction}\par\par\vspace{7mm}{\color{otRule!65}\hrule height .25pt}\par\vspace{7mm}{\color{otRule!65}\hrule height .25pt}\par\vspace{7mm}{\color{otRule!65}\hrule height .25pt}\par\vspace{7mm}{\color{otRule!65}\hrule height .25pt}\par\vspace{7mm}{\color{otRule!65}\hrule height .25pt}\par\vspace{7mm}{\color{otRule!65}\hrule height .25pt}
\end{blocEdition}
\editionPage{bilan-B-ex-2}{Bilan B · Problème 2}{Le marathon des enfants}{Partie 2 · 1 h 10 pour les quatre problèmes · 3,5 points}
\begin{blocEdition}{bilan-B-ex-2}
\begin{questions}
\item Une course démarre à 9 h 45 min 30 s. Le premier arrive après 38 min 45 s. À quelle heure arrive-t-il ?
\item Le dernier arrive à 10 h 41. Combien de temps a-t-il couru ?
\item Quel écart de temps sépare le premier du dernier ?
\end{questions}\par\vspace{5mm}{\otLight\fontsize{8}{10}\selectfont\color{otMuted}Recherche et rédaction}\par\par\vspace{7mm}{\color{otRule!65}\hrule height .25pt}\par\vspace{7mm}{\color{otRule!65}\hrule height .25pt}\par\vspace{7mm}{\color{otRule!65}\hrule height .25pt}\par\vspace{7mm}{\color{otRule!65}\hrule height .25pt}\par\vspace{7mm}{\color{otRule!65}\hrule height .25pt}\par\vspace{7mm}{\color{otRule!65}\hrule height .25pt}
\end{blocEdition}
\editionPage{bilan-B-ex-3}{Bilan B · Problème 3}{Le sondage de la classe}{Partie 2 · 1 h 10 pour les quatre problèmes · 3,5 points}
\begin{blocEdition}{bilan-B-ex-3}
\begin{questions}
\item Sur 28 élèves, la moitié vient à pied, le quart en bus, le reste à vélo. Donner les trois effectifs.
\item Dans un diagramme circulaire, quelle fraction du disque occupe le secteur « vélo » ?
\item On tire au hasard la fiche d’un élève. Quelles sont les chances de tirer un élève venu en bus ?
\end{questions}\par\vspace{5mm}{\otLight\fontsize{8}{10}\selectfont\color{otMuted}Recherche et rédaction}\par\par\vspace{7mm}{\color{otRule!65}\hrule height .25pt}\par\vspace{7mm}{\color{otRule!65}\hrule height .25pt}\par\vspace{7mm}{\color{otRule!65}\hrule height .25pt}\par\vspace{7mm}{\color{otRule!65}\hrule height .25pt}\par\vspace{7mm}{\color{otRule!65}\hrule height .25pt}\par\vspace{7mm}{\color{otRule!65}\hrule height .25pt}
\end{blocEdition}
\editionPage{bilan-B-ex-4}{Bilan B · Problème 4}{La figure à construire}{Partie 2 · 1 h 10 pour les quatre problèmes · 3,5 points}
\begin{blocEdition}{bilan-B-ex-4}
\begin{questions}
\item Construire un carré ABCD de côté 6 cm. Quel est son périmètre ? son aire ?
\item Tracer le cercle de centre A et de rayon 3 cm. Coupe-t-il le côté [AB] ? En quel point ?
\item Le carré a-t-il un axe de symétrie passant par A ? Lequel ?
\end{questions}\par\vspace{5mm}{\otLight\fontsize{8}{10}\selectfont\color{otMuted}Recherche et rédaction}\par\par\vspace{7mm}{\color{otRule!65}\hrule height .25pt}\par\vspace{7mm}{\color{otRule!65}\hrule height .25pt}\par\vspace{7mm}{\color{otRule!65}\hrule height .25pt}\par\vspace{7mm}{\color{otRule!65}\hrule height .25pt}\par\vspace{7mm}{\color{otRule!65}\hrule height .25pt}\par\vspace{7mm}{\color{otRule!65}\hrule height .25pt}
\end{blocEdition}
